\section{Three-dimensional storage with a protected classical calendar}\label{sec:quantum}

Write the physical coordinates as \((i,\sigma,z)\). Quantum data move
only along \(\sigma\), at fixed \((i,z)\). One semantic quantum qubit
per chain cell suffices: code coordinates and gadget work occupy more
cells. Two autonomous classical fields have distinct jobs. The selector
\(C_{\rm sel}\) supplies the compulsory final cut; a field \(S\)
schedules quantum correction. Every syndrome and every bit computed
from quantum data belongs to the quantum track, including when it is
basis-valued. This section constructs the positive storage circuit: a
one-level fault-tolerance bound (Lemmas~\ref{d:code}--\ref{d:reduction}),
a classical calendar and the source law for its errors
(Lemma~\ref{d:sources}), a renormalization bound (Lemma~\ref{d:recursion})
and storage at every requested time (Proposition~\ref{d:storage}).

The hierarchy follows the reliable-computation tradition of G\'acs
and G\'acs--Reif \cite{Gacs2001,GacsReif1988}; hierarchical local quantum
decoding has an early precedent in Harrington \cite{Harrington2004}.
We use the coherent CSS gadget designs of Aharonov and Ben-Or
\cite{ABO2008,ABO1999}. The support bound needed here is proved below
for the actual real-gate subsystem.

\subsection{An explicit binary code and one recovery}

\begin{lemma}[Binary code family]\label{d:code}
For every integer \(s\ge1\) there is an effective CSS code with parameters
\[
 [[n_s,1,d_s]]=[[2^{2s+1}-1,1,2^{s+1}-1]],
 \qquad q_s=2^s-1
\]
and one stabilizer recovery correcting the union span of all errors of
weight at most \(q_s\). Encoding, recovery and the total classical
functions used by the gadgets have finite exact real circuits.
\end{lemma}
\begin{proof}
Let \(\mathcal R_s=\mathrm{RM}(s,2s+1)\), the evaluations of multilinear
polynomials of degree at most \(s\) on \(\mathbb F_2^{2s+1}\).
Products of two generating monomials have degree below \(2s+1\),
so their evaluation inner product is even. Also
\(\dim\mathcal R_s=\sum_{j=0}^s\binom{2s+1}{j}=2^{2s}\), half the
length. Thus \(\mathcal R_s\) is self-dual. Each generator has weight
divisible by four. Self-orthogonality makes every pairwise intersection
even, and \(\mathrm{wt}(u+v)=\mathrm{wt}(u)+\mathrm{wt}(v)
-2|\supp u\cap\supp v|\) proves double evenness for every word.

The usual splitting \(f=f_0+x_m f_1\) proves the distance of
\(\mathrm{RM}(r,m)\) to be \(2^{m-r}\): if \(f_1=0\), the weight
is twice that of \(f_0\); otherwise the weights of \(f_0\) and
\(f_0+f_1\) sum to at least the weight of \(f_1\). Induction gives
the bound, and a product of \(r\) coordinate literals attains it.
Puncture at zero to obtain \(\mathcal C_s\). Its dual
\(\mathcal D_s=\mathcal C_s^\perp\) is the shortened code, contained
in \(\mathcal C_s\) with codimension one. The logical words are
\[
 |a_L\rangle=|\mathcal D_s|^{-1/2}
       \sum_{w\in\mathcal D_s}|w+a\one\rangle,
       \qquad a\in\mathbb F_2.
\]
Puncturing lowers distance by at most one, and
\(\prod_{j=1}^s(1+x_j)\) has value one at the removed point and
attains weight \(2^{s+1}-1\) in the nontrivial coset. This gives
\(d_s\) and \(q_s\). Products of two weight-\(q_s\) Pauli errors
have weight below \(d_s\); their compression to the code is scalar
or zero. The Knill--Laflamme equations therefore give one recovery
for the entire error span, including coherent cross terms.

Choose generators of \(\mathcal D_s\), complete them to an invertible
binary map, prepare the generator variables with Hadamards and the
remaining variables in zero, and apply that map by CNOT and SWAP.
This prepares the zero word; its transversal Hadamard gives the plus
word. Finite syndrome tables with a fixed choice outside the decoding
radius give total recovery circuits. Boolean functions use the
compute, copy and uncompute construction of
Section~\ref{sec:compilation}. All these prescriptions terminate for
each fixed \(s\). The code size and the running time may be enormous.
\end{proof}

\subsection{Real native procedures and the sign bit}
The calendar issues semantic gates, which Section~\ref{d:dispatch}
compiles into physical gates. Use the semantic inventory
\begin{equation}\label{d:native}
 \mathcal G_{\mathbb R}=
 \{I,X,H,\mathrm{CNOT},\mathrm{CZ},\mathrm{Toffoli},
                      \mathrm{SWAP},R_0,\text{discard}\}.
\end{equation}
Discard on an existing physical slot is reset. The transversal logical
identities follow directly from
\((a\one+w)\cdot(b\one+w')=ab\pmod2\) for
\(w,w'\in\mathcal D_s\), and character orthogonality for the Hadamard.
They give logical \(X,H,\mathrm{CNOT},\mathrm{CZ}\); SWAP exchanges
blocks. Toffoli is implemented coherently with a resource state below.

For a basis word \(w\in\mathbb F_2^{n_s}\), define \(a(w)\) by
bounded-distance decoding: if \(w\) is within Hamming distance \(q_s\)
of \(\mathcal C_s\), choose its unique nearest codeword and return
its coset in \(\mathcal C_s/\mathcal D_s\). Otherwise set \(a(w)=0\).
Uniqueness follows from \(d(\mathcal C_s)\ge2q_s+1\). In particular,
any at most \(q_s\) bit errors on a word in either logical coset
leave \(a(w)\) unchanged. Every decoded logical bit below, of a cat
block or of a consumed resource, data or target block, uses this decoder.

The sign operation at protected coordinate \(j\) is
\begin{equation}\label{d:sign}
 |w,b_j,c_j,d_j\rangle\longmapsto
 (-1)^{a(w)(b_jc_j+d_j)}|w,b_j,c_j,d_j\rangle.
\end{equation}
Compute \(a(w)\) in private scratch, apply
\(\mathrm{CCZ}(a,b_j,c_j)\mathrm{CZ}(a,d_j)\), then uncompute.
Here
\[
 Z=HXH,\quad \mathrm{CZ}=H_{\rm target}\mathrm{CNOT}H_{\rm target},
 \quad \mathrm{CCZ}=H_{\rm target}\mathrm{Toffoli}H_{\rm target}.
\]
The module touches one expendable cat block and only coordinate \(j\)
of each protected block. Its possibly large Boolean computation has
no other data target.

\begin{lemma}[Complete routed gadget bound]\label{d:gadgets}
The code family admits coherent error correction, deterministic verified
zero preparation, and coherent encoded bodies for
\(\mathcal G_{\mathbb R}\), with one-dimensional routing. A fault-free
complete error correction maps every input into the code. For a
correctable input deviating on at most \(f\) coordinates per block,
a complete routed error-correction-plus-body rectangle with \(j\)
faulty locations has output in the operator span of the ideal output
with at most \(8j\) erroneous coordinates per block, provided
\begin{equation}\label{d:margin}
 8(f+j)\le q_s-8.
\end{equation}
The statement includes arbitrary environment and reference coefficients.
Physical decomposition multiplicity is counted in the location budget,
rather than in the constant eight.
\end{lemma}
\begin{proof}
We give the four procedure arguments. A native location is a gate, wait
or reset of the circuit before routing. All records and their reversible
processing are quantum work. A private correction decision for coordinate
\(i\) controls only that coordinate. Record copying is done before
these private computations. Thresholds and repetitions below exceed
the admitted fault budget.

\smallskip\noindent
\emph{Deterministic verified preparation.}
Let \(Q_0\) be the exact encoder for \(|0_L\rangle\).
For \(t\ge1\), \(Q_t\) uses five \(Q_{t-1}\) blocks: a candidate
\(B_1\), testers \(B_2,B_3,B_4\), and spare \(B_5\).
Couple \(B_1\to B_2\) transversally and decode \(B_2\) modulo
\(\mathcal D_s\), including its logical bit. First check the
bit-error pattern of \(B_3\) with \(B_4\); then couple
\(B_3\to B_1\), apply \(H\) to \(B_3\) and decode the phase check.
Each coordinate reads private copies of the three decoded patterns.
It keeps \(B_{1,i}\) if their union has weight at most \(t\), and
otherwise Fredkin-swaps it with \(B_{5,i}\). Dispose of the work at
its declared end. This is a predetermined replacement circuit.

With \(k\le t\) faults, induction gives an output in the span of
\(F|0_L\rangle\) with \(\mathrm{wt}(F)\le k\).
If every recursive subtree has at most \(t-1\) faults, its deviation
is at most its own fault count. Uncorrupted private decisions see a
union bounded by the total count \(k\le t\), and keep the candidate.
Its residual consists of its own errors, phase errors transferred
from \(B_2\), bit errors transferred from \(B_3\), and coordinates
hit in the final verification/selection circuit. Each is charged to
its originating fault; their total is at most \(k\). A faulty private
computation can change only its own output coordinate.
If one subtree has at least \(t\) faults, it consumes the entire
budget: \(k=t\) and the final checks and all other blocks are clean.
All private decisions then agree. A bad candidate is either replaced
by the clean spare or accepted with union pattern of weight at most
\(t\). A bad \(B_2\)'s transferred phase errors are detected by the
exact phase check; a bad \(B_3\)'s dangerous bit pattern is checked
by \(B_4\) before coupling. If only \(B_4\) is bad, rejection selects
the clean spare; if only \(B_5\) is bad, the clean candidate passes.
These cases give the same bound. In particular the output cannot be
a coordinatewise splice of two arbitrary codewords: in the first
case every uncorrupted decision keeps, and in the second all
decisions agree. Use depth \(t=q_s\); it is finite and at least the
native rectangle budget \(2b_{\rm FT}\) selected below.

\smallskip\noindent
\emph{Coherent error correction.}
For bit correction couple data transversally into a verified plus block.
For each eventual data coordinate make a private record and compute
its nearest-codeword error vector modulo \(\mathcal C_s\); apply
only that vector's corresponding correction bit. Phase correction
is the same operation in the Hadamard basis. A data-record copying
fault changes one coordinate in each record. Within the unique-decoding
margin it changes the inferred error by that same coordinate; a fault
inside a private computation affects only its own data target.
The transversal coupling transmits at most one coordinate from a
verified-ancilla error. The preparation bound just proved charges
those coordinates to their actual faults. Applying the two stages
removes the incoming correctable error and leaves at most one
coordinate per native fault in each output block.

For normalization, also compute the radius indicator privately.
In an exact syndrome branch outside the prescribed decoding region,
exchange every data coordinate with the corresponding coordinate of
a verified zero block. In a branch inside it, correct the error.
In the fault-free circuit all private indicators agree in each branch,
and both branches have code support. Thus the entire output has code
support for an arbitrary input, including a purification. Under the
margin used here the intended correctable branch is separated from
the replacement threshold; a disturbed private decision has only its
own coordinate as target. There is no shared faulty accept bit.

\smallskip\noindent
\emph{Coherent Toffoli.}
Put
\[
 |A_0\rangle=\frac12\sum_{a,b}|a_L,b_L,(ab)_L\rangle,
 \qquad |A_1\rangle=X_{L,3}|A_0\rangle.
\]
Their equal superposition is three encoded plus words. Repeated encoded
cat sign tests using \eqref{d:sign} distinguish these branches.
Use \(R_{\rm cat}=2\lfloor q_s/2\rfloor+1\) tests, complete EC on
the resource blocks between tests, Hadamards on the cats, and private
per-coordinate computations that decode each cat block by the
bounded-distance coset decoder, take the parity of those decoded bits,
and compute the majority over tests. A private
majority controls only its coordinate of the branch-correcting \(X\).

One native fault spoils at most one repeated test and touches at most
one coordinate of each protected block. A residual resource-coordinate
error is removed by the next EC. During cat fan-out, an error on one
coordinate of the first cat block can be copied to that coordinate
of later cat blocks. Each such block remains in the same decoded
coset, so \(a(w)\) is unchanged. A fault in one sign computation can
spoil its expendable cat and its one protected coordinate; that cat
is not reused on other coordinate triples. Thus cat fan-out does not
turn one fault into many wrong signs on protected coordinates.
The majority is correct while the native fault count is at most
\(\lfloor q_s/2\rfloor\). Private parity/majority faults have the
same one-target bound. This proves the resource support assertion.

The resource-to-gate conversion uses transversal CNOT and Hadamard,
private decoding of consumed blocks, and private real corrections.
Its disposal preserves an unknown complex qubit. For half-Toffoli on
logical data bits \(c,d\), let the resource carry \((x,y,xy)\); CNOTs
from its first two blocks into the data leave consumed outcomes
\(u=c+x\), \(v=d+y\). Add
\(xv+yu+uv\) to the third register and apply \(X^u,X^v\) to
the first two. Since \(xy+xv+yu+uv=(x+u)(y+v)=cd\), the
retained registers become \((c,d,cd)\) in every branch;
\((u,v)\) is uniform independently of \((c,d)\), and all decoder
garbage depends on \((u,v)\) alone. For full Toffoli on data bits
\((a,b,c)\) the new target is \(c+ab\). The old target, after Hadamard,
decodes to a bit \(h\) and contributes phase \((-1)^{hc}\); the output correction
contributes
\[
 (-1)^{h(ab+(c+ab))}=(-1)^{hc}.
\]
They cancel, leaving the same garbage vector for every \((a,b,c)\).
Consequently disposal gives the correct gate for arbitrary complex
amplitudes and a reference. Each conversion fault has one protected
coordinate per block as its contact. The complete EC followed by
this body, as well as each transversal body, first removes the
correctable input and obeys the same total-count support bound.

\smallskip\noindent
\emph{Routing and the numerical envelope.}
Route the already complete EC-plus-body circuit along a line. Bring
its at most three operand tokens together by a fixed adjacent-SWAP
word, perform the native gate, and undo the permutation. Include
identity locations on waiting tokens. Reset a work lifetime before
its first use and retain it until its declared reset or disposal.
An ideal SWAP permutes errors. Expand a faulty two-token SWAP in a
product operator basis and commute its factors through the remaining
permutations. The operand factor is an arbitrary fault at the next
native gate; a bypass token's factor is a fault at its native idle.
Return routes have the same description. Thus one routed fault gives
at most two native faults, with possibly correlated coefficients,
independently of route length and \(s\). Waits and resets give one.

The premise \eqref{d:margin} implies
\[
 f+2(j+1)\le2(f+j+1)\le q_s/4\le\lfloor q_s/2\rfloor.
\]
All unique-decoding and cat-majority arguments are therefore inside
their stated margins, including one additional insertion. The routed
total deviation is at most \(2j\), hence at most \(8j\).
Route length, private table size and the \(5^{q_s}\) preparation
leaves affect location count and time, not contact per fault.

Finally these are operator statements. Expand arbitrary faulty Kraus
operators in local matrix or Pauli bases, retaining all exterior
systems as coefficients. In EC a coset-uniform record
\(\sum_{v\in\mathcal C_s}|v+x\rangle\) depends on the error coset,
not the logical amplitude. The two Toffoli cancellations show the
same independence for its consumed records. On the correctable span
each term is an ideal logical map times an input-independent
syndrome/environment coefficient. The support analysis applies
coefficientwise, so coherent sums, cross terms and the full reference
are covered by the one recovery of Lemma~\ref{d:code}.
\end{proof}

All procedures in this proof use \(\mathcal G_{\mathbb R}\).
An encoded Toffoli table contains physical Toffolis; replacing them
at the next level is finite circuit substitution. Induction on the
number of levels proves that every finite concatenated identity
circuit remains in this inventory. Invalid opcodes select identity.
Section~\ref{d:dispatch} reduces it to the seven physical gates.

\subsection{Hardened correction and correlated normalizers}
Fix the constant semantic width, candidate-port geometry and complete
boundary halo before selecting a code. Let \(\kappa_0\) count physical
conversion into dirty semantic locations, \(\kappa_1\) bound child-box
participation, and \(v_0\) bound parent boxes meeting a rectangle.
These are incidence bounds per cell and time; extra ancillas use more
cells and extra correction passes use more time. With complete halo
\(h\) and a fixed shadow gap \(\rho\), choose
\begin{equation}\label{d:budgets}
 a_{\rm sep}>(4h+2\lceil\rho\rceil+5)^2,
 \qquad b_{\rm FT}=\max\{1,\kappa_0\kappa_1v_0(a_{\rm sep}-1)\},
\end{equation}
and then choose
\begin{equation}\label{d:capacity}
 e=8,\quad r_{\rm FT}=8b_{\rm FT},\quad
 s=\min\{s\ge1:2^s-1\ge40960(b_{\rm FT}+1)\}.
\end{equation}
All Fredkin, CZ, SWAP, CCZ and Z decomposition faults are paid in
\(\kappa_0\), hence \(b_{\rm FT}\). The preparation depth is
\(q_s\ge2b_{\rm FT}\). The integer \(s\) is fixed for the medium.

Let \(B_g\) be one full EC followed by an
encoded body. Precede every \(B_g\) by \(b_{\rm FT}+1\) time-disjoint
full EC passes on each input block, with fresh private work. Synchronize
these normalizing phases before interblock bodies. Every idle block
normalizes too.

\begin{lemma}[Joint level reduction]\label{d:reduction}
If a parent rectangle contains at most \(b_{\rm FT}\) bad child
circuit locations, its output is correctable by one fixed recursive
recovery. Good logical subsystems are preserved jointly, including
when entangled with bad blocks and an arbitrary reference. Effective
faults at the next level have support only on the union of the named
bad child locations' outgoing wires.
\end{lemma}
\begin{proof}
Each good block's normalization has a last clean EC pass, because
there are more full passes than touching faults. That pass maps any
input into the code. Every later full EC pass and \(B_g\) has input
deviation at most \(r_{\rm FT}\) and at most \(b_{\rm FT}\) faulty
locations, and
\[
 8(r_{\rm FT}+b_{\rm FT})=72b_{\rm FT}\le q_s-8.
\]
Lemma~\ref{d:gadgets}, applied afresh at each full EC and at \(B_g\),
leaves at most \(r_{\rm FT}\) residual coordinates. A correctable
input retains its original logical subsystem. An arbitrary input has
output support in
\[
 \mathcal C_{r}=\operatorname{span}
       \{FV|\psi\rangle:\mathrm{wt}(F)\le r_{\rm FT}\}.
\]
Different good blocks can have their last clean pass at different times.

Actual faults can couple blocks. For each good block expand shared
operators on its coordinates, with all other systems as coefficients.
The preceding argument puts each output in \(\mathcal C_r\).
Positivity then puts the joint output in the tensor product of these
spaces: every complementary projection has expectation zero. In fixed
recovery coordinates \(\mathcal C_r\simeq\mathcal H_{\rm log}
\otimes\mathcal H_{\rm syn}\), the complete operator identity in
Lemma~\ref{d:gadgets} fixes the whole originally good logical algebra.
Its Stinespring map is
\[
 I_{\text{all good logicals}}\otimes
       K_{\text{bad inputs, syndromes, memory}}.
\]
The coefficient is the same for every good logical amplitude, including
off-diagonals; the factor on the remaining systems may be correlated.

At a common coarse cut, compose the next normalizations with the
preceding bad outgoing wires before crossing any interblock body.
In recovery coordinates their joint action is identity on the good
logical wires and a joint map on the union of those bad wires and
syndrome memory. Expand that map in tensor matrix bases on the named
bad wires. It adds neither a wire nor a bad location. Consecutive bad
rectangles can remain distinct; their common environment coefficients
are allowed. Good bodies now implement their ideal logical gates.
The finite directed circuit permits this reduction successively at
every level. Full scheduled normalizers remain covered at their own
scale; their circuit is never truncated to one fine time step.

For completeness, the decoder is fixed bottom-up. At a parent there
are at most \(r_{\rm FT}\) arbitrary child-block operators; all other
children are recursively correctable. Lower recovery turns the former
into arbitrary operators on at most \(r_{\rm FT}\) parent coordinates,
with exterior coefficients. The base recovery corrects their entire
union span. Induction corrects all recursively allowed supports and
their coherent combinations. The same decoder serves every history.
\end{proof}

\subsection{A finite classical interpreter for the schedule}\label{d:calendar-section}
Choose a sufficiently large even block length \(P_{\rm cal}\) and a
period \(T_{\rm cal}=P_{\rm cal}^4\), or one larger fixed polynomial.
The radius-one classical field \(S\) has finite address and age,
a mode, distributed old/next information bits, read-only program bits,
and finite tape/head and broadcast tracks. A block of
\(P_{\rm cal}\) cells stores one virtual state as a binary string of
length \(O(\log P_{\rm cal}+\log T_{\rm cal})\), with blank padding.

During a simulation period it copies its old string and those of its
two neighbors, interprets one local transition and its old quantum
instruction, broadcasts the instruction, executes its padded encoded
quantum table, and commits the next string. Communication and broadcast
cost \(O(P_{\rm cal})\); one virtual transition uses polynomial time
and space in the logarithms of the parameters. The algorithm computes
one transition, rather than running the entire virtual period.
For the already fixed quantum table, padding and line routing take
at most \(C_QP_{\rm cal}^2\) steps. Include quantum idles during
classical work, the hardened leading EC, and every reset. Thus a
terminating search finds \(P_{\rm cal}\) large enough for all work
to fit. Pair instructions use fixed parity/matching roles, with even
parameters allowing the same rule recursively. Deadlines and fixed
conflict priorities make malformed classical words total.

The local algorithm is defined first. The encoder writes its finite
description as ordinary program data. The interpreter computes that
algorithm on three virtual states, including their program-bit fields;
it does not require a self-description fixed point. In the idle mode
the virtual instruction is identity. A simulation block has the exact
identities
\[
 \begin{aligned}
 \text{new virtual state}&=\text{one classical transition},\\
 \text{quantum channel}&=E_{\mathbb R}(\text{old virtual instruction}).
 \end{aligned}
\]
Induction gives a finite tower of \(k\) simulation levels under a top
idle state, executing \(E_{\mathbb R}^k(I)\). Higher clocks and
programs stay classical. Every finite initial symbol is computed by
\(k\) block decompositions and finite encodings.

\subsection{Protecting the entire calendar symbol and counting sources}
Clone the intended \(S\) trajectory over the transverse \((i,z)\)
plane. At every physical update predict its next whole symbol from
its three old \(\sigma\)-neighbors at every provider in
\eqref{eq:c-boxes}; retain a symbol with at least 512 votes, or a
fixed default. A clean wrong output needs at least 257 wrong
predictions, occupying two provider boxes. Each wrong prediction
has a wrong old input. Use private copies of old \(S\) for the
current quantum instruction, before committing its next symbol.
A later correct retained symbol does not undo a wrong current use.

\begin{lemma}[Combined source certificates]\label{d:sources}
Every primitive dirty quantum location has a connected full witness
with \(g\) distinct canonical source labels and at most \(C_{\rm w}g\)
edges. Its anchored description count is at most \(K_{\rm w}^g\).
Under the one joint physical law, any requested \(g\) labels have
intersection probability at most \(\gamma_{\rm eff}^g\), where
\[
 \gamma_{\rm eff}=\max\{\sqrt{\gamma_0},\sqrt p\}.
\]
Each source has one Lipschitz shadow in chain space-time. Full bin
and role information is retained when witnesses are counted.
\end{lemma}
\begin{proof}
For \(S\) use characteristics \(i+58t,z+58t\) and balanced poles
\(-3(i+58t),-3(z+58t),3(i+z+116t)\). The three provider boxes give
ranges
\[
\begin{array}{c|rrr}
 &F_1&F_2&F_3\\ \hline
 S_0 &[129,174]&[129,174]&[-348,-258]\\
 S_1 &[-351,-306]&[129,174]&[132,222]\\
 S_2 &[129,174]&[-351,-306]&[132,222]
\end{array}
\]
for backward gains. Any two boxes provide a positive parent for each
pole, with margin at least 129. The \(\sigma\)-shift affects none
of these poles. Lemma~\ref{q:frontier}'s grouped-terminal proof
therefore gives \(O(g_H)\) edges for distinct actual raw names
\((i,\sigma,z,t)\). Old/new roles and quantum hardware use the same
raw hub \(H_{i,\sigma,z,t}\). The literal shadow \((\sigma,t)\)
is Lipschitz; there are no helper contractions in this raw graph.

For \(C_{\rm sel}\), the canonical hub is
\[
 \lambda=(i-t,\sigma+58t,z+58t,\lfloor t/W\rfloor).
\]
Its event is the union of its raw, statistical and overlay exceptions,
with the law \(\gamma_0^{g_C}\) in Proposition~\ref{c:source-law}.
These three types annotate one hub. Helper gather/consult links cancel
on the base characteristic; a clean predictor entry traverses all
1024 levels, distributing its fixed debit over those levels.
The selector link calculation and grouped-terminal argument are given
in Table~\ref{q:links} and Lemma~\ref{q:frontier}; deleting the copied
coordinate changes neither. Keep the bin and entry roles in the full
bounded-degree graph. On a quantum line \(i=i_0\) its shadow is
\[
 t^*=i_0-(i-t),\qquad
 \sigma^*=(\sigma+58t)-58t^*.
\]
It agrees with the actual root and is constant on each canonical hub.
Contracted helper links have zero shadow increment; all other typed
links have fixed bounded increments.

Glue the two full graphs at actual quantum-use roots by the fixed
hardware stars. For \(g_C\) owner hubs and \(g_H\) raw hubs,
Cauchy--Schwarz gives
\[
 \Pr(\text{all})\le\gamma_0^{g_C/2}p^{g_H/2}
                  \le\gamma_{\rm eff}^{g_C+g_H}.
\]
This pays every physical overlap. Bounded full degree and the linear
edge budget give exponential anchored enumeration by a doubled-tree
walk with finite annotations. In particular neither source types nor
coincident shadows create duplicate charges or free bin choices.
\end{proof}

\subsection{Separated complete supports and the threshold}\label{d:threshold-section}
Use regular chain space-time boxes of spatial side
\(P_{\rm cal}^{\ell}\) and temporal side \(T_{\rm cal}^{\ell}\).
Their fixed child-index cover includes a constant collar. Geometric
series bound all complete descendants by a fixed halo \(h\).
Let \(N_0\) mean primitive dirtiness, and define \(N_{\ell+1}\)
by at least \(a_{\rm sep}\) bad children in that cover. This is a
monotone geometric predicate on scheduled bad roots. If all parent
boxes covering a rectangle are good, \eqref{d:budgets} bounds its
bad child circuit locations by \(b_{\rm FT}\). Thus
Lemma~\ref{d:reduction} applies.

The packing bound in \eqref{d:budgets} ensures two bad complete child
supports separated by \(\rho P_{\rm cal}^{\ell}\) whenever the
parent is bad. Separation in time can be still larger, because
\(T_{\rm cal}^{\ell}\ge P_{\rm cal}^{\ell}\). There are finitely
many candidate pairs per parent, independently of \(\ell\).

\begin{lemma}[Source-certificate recursion]\label{d:recursion}
At a sufficiently small computably selected source cap there are
constants \(\bar A\ge1\) and \(D=1/(2\bar A)\) such that
\[
 \Pr(N_\ell)\le z_\ell\le D\,2^{-2^\ell}
\]
uniformly in the hierarchy level and selected nonhalting window.
\end{lemma}
\begin{proof}
Let \(z_\ell\) be the sum of source weights of recursively pruned
fully embedded certificates in one level-\(\ell\) box. At a parent,
if a selected primitive witness has full graph length at least a
small fixed multiple of \(P_{\rm cal}^{\ell}\), replace its entire
repeated subtree by that witness and its root. Lemma~\ref{d:sources}
pays at most
\(B_{\rm tail}P_{\rm cal}^{d_0\ell}
 e^{-\alpha P_{\rm cal}^{\ell}}\).
The polynomial counts possible roots, including the anisotropic time
extent. Otherwise every witness is shorter than one third of the
separated-support gap. The two complete child certificates have
disjoint source-label sets, since a label has one shadow. Their
weights multiply under the joint intersection bound, and enumeration
factors. Hence
\begin{equation}\label{d:sum-recursion}
 z_{\ell+1}\le A_{\rm rec}z_\ell^2+
 B_{\rm tail}P_{\rm cal}^{d_0\ell}e^{-\alpha P_{\rm cal}^{\ell}}.
\end{equation}
The long case discards the full subtree; its leaves are not counted
again. This is a recursion for certificate sums, not for independent
bad-location probabilities.

Put \(\bar A=\max(1,A_{\rm rec})\). Choose the source cap so that
\[
 z_0\le\frac1{4\bar A},\qquad
 \alpha\ge\max\{0,\log(4\bar A B_{\rm tail})\}
       +\frac{d_0\log P_{\rm cal}}{P_{\rm cal}}+2\log2.
\]
For \(P_{\rm cal}\ge3\), use
\(\ell/P_{\rm cal}^{\ell}\le1/P_{\rm cal}\) for \(\ell\ge1\)
and \(2(2/P_{\rm cal})^\ell\le2\).
The remainder in \eqref{d:sum-recursion} is at most
\((D/2)2^{-2^{\ell+1}}\). If \(z_\ell\le D2^{-2^\ell}\), its
squared contribution is at most the same quantity because
\(A_{\rm rec}D^2\le D/2\). The base inequality is
\(z_0\le D/2\). This proves the displayed bound with the fixed
exponent \(\log2\). All entropy constants precede the source cap;
shrinking that cap makes these strict conditions attainable.
\end{proof}

\subsection{A total encoder and every requested time}

\begin{proposition}[Three-dimensional linear storage]\label{d:storage}
There are a rational \(c>0\), independent of \(n\), and total
computable \(L_0(n)\), \(\cE_{n,L}\), \(\cD_{n,L,t}\) satisfying
part (i) of Theorem~\ref{thm:detailed}. One encoder works for every
requested time and every rate in its window.
\end{proposition}
\begin{proof}
Let \(k\) be the least integer \(k\ge1\) with
\(2^k\ge64\lceil\log_2(L+2)\rceil\), and set
\[
 t_{\max}=\lfloor cL\rfloor,\qquad
 t_{\rm ext}=t_{\max}+T_{\rm cal}^{k}.
\]
The hierarchy dimensions \(P_{\rm cal}^k,T_{\rm cal}^k\) are fixed
powers of \(\log(L+2)\), hence \(o(L)\).
Choose \(0<c\le1/[1000(V_*+1)]\) after the final catalogue and full
source speed \(V_*\) are known. A total computable size threshold
makes the complete extended quantum and source cones inject into
one torus chart, with all seams outside. Fixed label-dependent
initialization and bin margins enter \(L_0(n)\).

Prepare the finite \(k\)-level quantum block and classical simulation
tower under an idle top. Clone intended \(S\) over the required
transverse cone. Lemma~\ref{c:initialization} supplies the finite
selector chart by exhaustive hidden-cover search, with low/high
helpers on their base rows. Nonhalting guarantees success. The search
is finite even for halting inputs; failure and small sizes use a fixed
fallback. Neither a noise rate nor a halting answer is read.

At most \(C_*L^4\) complete level-\(k\) boxes meet the extended
cone and its fixed cover. Define \(\mathcal G\) by the absence of
\(N_k\) in all of them. Lemma~\ref{d:recursion} gives
\[
 \Pr(\mathcal G^c)\le C_*L^4D(L+2)^{-64}\le1/16
\]
after a computable threshold enlargement. The event is exogenous
and independent of the stored state. On it the deterministic
reduction restores the ideal top identity on its full logical algebra.

At a requested integer \(t\le t_{\max}\), retain every actual
in-flight quantum data, syndrome, partial correction, ancilla and
routing/work factor. Ignore the received autonomous metadata and
append the intended noiseless future quantum circuit until
\(t'=T_{\rm cal}^k\lceil t/T_{\rm cal}^k\rceil\).
Existing work is reset only at a declared future reset instruction;
new external work is introduced only as genuinely fresh work.
Physical disposable ports already cleaned at the round boundary may
be discarded. This external completion adds neither noise nor
selector cuts. Delete only future primitive bad roots in the comparison.
Monotonicity of the pure \(N_\ell\) predicates preserves goodness,
so Lemma~\ref{d:reduction}'s one bottom-up recovery gives the Bell
state exactly. Averaging gives fidelity at least \(15/16\).
The initial encoder did not depend on the requested time. The
external circuit is total and computable from \((n,L,t)\), independently
of the actual rate and fault history.
\end{proof}
